\documentclass[11pt, a4paper]{article}

% --- Codificação e Idioma ---
\usepackage[utf8]{inputenc}
\usepackage[T1]{fontenc}
\usepackage[brazil]{babel}

% --- Fonte e Tipografia (Estilo Matemático) ---
\usepackage{mathpazo} % Fonte Palatino (Texto e Matemática)
\usepackage[scaled=0.95]{helvet} % Helvetica para Sans Serif
\usepackage{courier} % Courier para código
\usepackage{microtype} % Micro-tipografia refinada
\linespread{1.05} % Espaçamento de linha Palatino

% --- Pacotes Matemáticos ---
\usepackage{amsmath}
\usepackage{amsfonts}
\usepackage{amssymb}
\usepackage{amsthm}
\usepackage{mathtools}

% --- Estilo dos Títulos ---
\usepackage{titlesec}
\titleformat{\section}{\large\scshape\centering}{\thesection.}{1em}{}
\titleformat{\subsection}{\bfseries}{\thesubsection}{1em}{}

% --- Legendas e Floats ---
\usepackage[font=small,labelfont=bf]{caption}
\usepackage{float}

% --- Algoritmos ---
\usepackage{algorithm}
\usepackage{algpseudocode}
\floatname{algorithm}{Algoritmo}
\renewcommand{\algorithmicrequire}{\textbf{Entrada:}}
\renewcommand{\algorithmicensure}{\textbf{Saída:}}

% --- Layout e Margens ---
\usepackage{geometry}
\geometry{
 a4paper,
 left=20mm,
 right=20mm,
 top=25mm,
 bottom=25mm,
}

% --- Cabeçalhos e Rodapés ---
\usepackage{fancyhdr}
\pagestyle{fancy}
\fancyhf{} 
\fancyhead[CO]{\small \scshape Projeto Vaccai}
\fancyhead[CE]{\small \scshape Documentação Técnica}
\fancyhead[R]{\small \thepage} 
\renewcommand{\headrulewidth}{0.5pt}

% --- Configuração de Teoremas ---
\theoremstyle{plain}
\newtheorem{teorema}{Teorema}[section]
\theoremstyle{definition}
\newtheorem{definicao}{Definição}[section]

% --- Links e Referências ---
\usepackage{xcolor}
\definecolor{azullink}{RGB}{0, 0, 100}
\usepackage{hyperref}
\hypersetup{
    colorlinks=true,
    linkcolor=azullink,
    citecolor=azullink,
    urlcolor=azullink,
    pdftitle={Projeto Vaccai - Documentação},
    pdfauthor={Rumenick Brandi Sigolo}
}

% --- Metadados do Título ---
\title{\textbf{\LARGE Projeto Vaccai}\\[0.2cm] \large Documentação Técnica e Modelagem Matemática}
\author{\textsc{Rumenick Brandi Sigolo}}

\begin{document}

\maketitle

% Resumo
\begin{center}
\begin{minipage}{0.9 \textwidth}
\begin{abstract}
\noindent Este documento centraliza toda a parte matemática e os algoritmos utilizados no motor de processamento de áudio (DSP) do Projeto Vaccai. O foco principal é a análise de voz monofônica em tempo real.

O texto passa por todo o processo atual: desde o tratamento inicial do sinal e a detecção de frequência fundamental com o algoritmo YIN. Além disso, aborda a quantização para MIDI e o uso de Modelos Ocultos de Markov (HMM) para identificar a tonalidade e o modo da música. Sempre que novas lógicas ou funcionalidades forem adicionadas ao projeto, este documento também será atualizado para refletir a nova modelagem.
\end{abstract}
\end{minipage}
\end{center}

\vspace{0.5cm}

\section{Introdução}
O Projeto Vaccai é um sistema de processamento de áudio em tempo real voltado para o treinamento vocal. O principal objetivo do motor de DSP (\textit{Digital Signal Processing}) é capturar a frequência fundamental ($f_0$) da voz com alta precisão e praticamente sem atraso (baixa latência), convertendo tudo para notas musicais para analisar a afinação e a tonalidade.

Esta documentação funciona como a base teórica e matemática oficial do \textit{backend}. À medida que o código evoluir e novas soluções forem integradas, a modelagem por trás delas será continuamente adicionada aqui.

\subsection{Justificativa da Abordagem no Domínio do Tempo}

Para a detecção de $f_0$, a escolha foi usar algoritmos no domínio do tempo (autocorrelação via YIN) em vez das clássicas abordagens espectrais via FFT. O motivo disso é o \textit{trade-off} fundamental entre resolução espectral e temporal, limitado pelo princípio de Gabor-Heisenberg\footnote{O Princípio da Incerteza de Gabor estabelece que não dá para localizar um sinal com precisão perfeita no tempo e na frequência ao mesmo tempo. O produto das incertezas no tempo ($\Delta t$) e na frequência ($\Delta f$) tem um limite inferior: $\Delta t \cdot \Delta f \ge C$, onde a constante $C$ depende de como se define a largura de banda.}.

Em métodos de transformada discreta, a resolução de frequência $\Delta f$ (a largura do \textit{bin}) é inversamente proporcional ao tamanho da janela de tempo $T_w$:

\begin{equation}
    \Delta f = \frac{f_s}{N} = \frac{1}{T_w}
\end{equation}

Onde $f_s$ é a taxa de amostragem e $N$ o número de amostras. No caso do canto, a precisão nas frequências graves é crítica. Por exemplo, a diferença entre as notas $A_2$ ($110$ Hz) e $A\#_2$ ($116.5$ Hz) exige uma separação de apenas $\approx 6.5$ Hz. Para conseguir resolver essa diferença direto na FFT, o tamanho mínimo da janela teria que ser:

\begin{equation}
    T_w \approx \frac{1}{6.5 \text{ Hz}} \approx 153 \text{ ms}
\end{equation}

Uma latência de $153$ ms, somada ao tempo de computação e ao \textit{buffer} de áudio, criaria um atraso perceptível. Isso inviabilizaria o \textit{feedback} visual instantâneo, que é o objetivo do Projeto Vaccai.

Por outro lado, métodos no domínio do tempo funcionam medindo a periodicidade do sinal (atraso $\tau$). Isso permite atingir uma precisão de frações de Hertz usando janelas bem menores (geralmente 2 a 3 períodos da fundamental), o que garante a resposta rápida do sistema.

\section{Modelo de Sinal e Pré-Processamento}

Para manter a baixa latência justificada na seção anterior, o sistema trabalha diretamente no fluxo de áudio bruto. A ideia é evitar o uso de janelas de acumulação muito extensas, o que criaria um atraso perceptível. Contudo, a eficiência dos algoritmos no domínio do tempo depende muito da qualidade do sinal de entrada.

Considerando $x[n]$ o sinal de áudio discreto com taxa de amostragem $f_s$ (amostras por segundo), onde $n$ é o índice da amostra no tempo, o processamento ocorre em janelas (\textit{buffers}) de tamanho fixo $N$. Esse tamanho é dimensionado justamente para garantir a resposta em tempo real.

Na análise vocal, o sinal $x[n]$ não é contínuo; ele é intercalado por pausas de silêncio e ruídos não-vocais (como respiração, som ambiente ou consoantes fricativas). O pré-processamento tem o papel crítico de ``limpar'' esse sinal e decidir, de forma determinística, quando o algoritmo de detecção de \textit{pitch} deve atuar. Isso evita o processamento de ruídos que gerariam resultados musicais espúrios (notas erradas). A lógica completa desse condicionamento e decisão de fluxo está detalhada no \textbf{Algoritmo \ref{alg:preprocess}}, no Apêndice A.

\subsection{Remoção de Componente Contínua (DC Offset)}
Sinais de voz capturados por microfones comuns costumam ter um deslocamento de tensão vertical, conhecido como \textit{DC Bias}. Isso acontece por conta de pequenas imperfeições no hardware de conversão analógico-digital, fazendo com que a onda não oscile exatamente em torno do zero.

Esse deslocamento atrapalha a análise de periodicidade, já que altera os valores no cálculo dos produtos $x[n] \cdot x[n+\tau]$. A solução para corrigir isso é centralizar o sinal subtraindo a sua média:

\begin{equation}
    \mu_x = \frac{1}{N} \sum_{n=0}^{N-1} x[n]
\end{equation}

\begin{equation}
    x'[n] = x[n] - \mu_x
\end{equation}

Essa operação traz a simetria da onda de volta, o que é fundamental para conseguir extrair o \textit{pitch} com precisão.

\subsection{Discriminação Híbrida de Voz e Silêncio (Histerese Dupla)}

A detecção de atividade vocal combina a energia do sinal (RMS) com a pureza harmônica para decidir de forma assertiva sobre a presença de voz. Para medir a energia, calcula-se a amplitude \textit{Root Mean Square} para cada janela:

\begin{equation}
    A_{rms}[t] = \sqrt{\frac{1}{N} \sum_{n=0}^{N-1} (x'[n])^2}
\end{equation}

Uma melhoria importante é a capacidade de diferenciar vogais de ruídos (como /s/ ou respiração) usando a Taxa de Cruzamento por Zero (ZCR):

\begin{equation}
    ZCR[t] = \frac{1}{N-1} \sum_{n=1}^{N-1} \mathbb{I}(x'[n] \cdot x'[n-1] < 0)
\end{equation}

Para decidir o estado de processamento $S_t \in \{0, 1\}$ (onde 1 é "Voz" e 0 é "Silêncio"), um comparador simples oscilaria rapidamente gerando \textit{chattering}. O sistema resolve isso implementando uma \textbf{histerese dupla}, que avalia simultaneamente amplitude e ZCR:

\begin{equation}
    S_t = \begin{cases} 
    1 & \text{se } A_{rms}[t] > \eta_{open} \land ZCR[t] < \zeta_{open} \\
    0 & \text{se } A_{rms}[t] < \eta_{close} \lor ZCR[t] > \zeta_{close} \\
    S_{t-1} & \text{caso contrário}
    \end{cases}
\end{equation}

Os parâmetros foram rigorosamente calibrados para garantir responsividade:
\begin{itemize}
    \item \textbf{Limiares de Abertura ($\eta_{open} = 0.003$, $\zeta_{open} = 0.20$):} O áudio deve ter volume suficiente e cruzar o zero poucas vezes (indicando um som vozeado limpo) para iniciar a análise.
    \item \textbf{Limiares de Fechamento ($\eta_{close} = 0.001$, $\zeta_{close} = 0.30$):} Se o áudio cair abaixo do piso de ruído ou se tornar excessivamente sibilante (ZCR alto), o motor de processamento entra em pausa.
\end{itemize}

Essa "zona morta" funciona como uma inércia digital, garantindo que o \textit{feedback} visual tenha brechas limpas e não trema durante finais de frase orgânicos.

\section{Detecção de Frequência Fundamental (Algoritmo YIN)}

Quando o pré-processamento confirma que existe atividade vocal estável ($S_t = 1$), o sistema avança para a extração da frequência fundamental $f_0$. Para essa tarefa, utiliza-se o \textbf{Algoritmo YIN} \cite{yin2002}, escolhido por ser significativamente mais robusto que métodos simples de contagem de cruzamentos por zero ou correlação padrão. A estrutura completa dessa implementação encontra-se no \textbf{Algoritmo \ref{alg:pitch}}, no Apêndice A.

\subsection{Função de Diferença e Normalização Cumulativa}
O núcleo matemático do YIN opera de forma inversa à correlação: em vez de buscar a máxima semelhança, ele busca minimizar a diferença entre o sinal condicionado $x'[n]$ e suas versões atrasadas. A função de diferença quadrática é definida como:

\begin{equation}
    d_t(\tau) = \sum_{j=1}^{W} (x'[j] - x'[j+\tau])^2
\end{equation}

Para evitar o erro de "atraso zero" e impedir que o algoritmo escolha harmônicos superiores por engano, aplica-se a Normalização pela Média Cumulativa (CMNDF). Essa operação cria a função $d'_t(\tau)$, que penaliza vales profundos em atrasos curtos caso eles não se sustentem na média, reduzindo drasticamente os erros de oitava:

\begin{equation}
    d'_t(\tau) = \begin{cases} 
    1 & \text{se } \tau = 0 \\ 
    d_t(\tau) \cdot \left[ \frac{1}{\tau} \sum_{j=1}^{\tau} d_t(j) \right]^{-1} & \text{se } \tau > 0 
    \end{cases}
\end{equation}

\subsection{Seleção de Limiar (Busca Híbrida)}

A literatura original sugere um limiar absoluto e rigoroso em torno de 0.10 para evitar erros de oitava definitivamente. No entanto, para fins práticos e pedagógicos, o motor DSP do Projeto Vaccai implementa uma estratégia de \textbf{busca híbrida}. 

Primeiro, o algoritmo procura o primeiro vale que cruze um limite de tolerância estrito de $\theta = 0.15$. Se a voz do estudante iniciante for muito "soprosa" (com muito ar vazando), a aperiodicidade será alta e esse limiar rigoroso pode não ser alcançado. Neste caso específico, o algoritmo realiza um \textit{fallback} para o mínimo global encontrado na janela de atrasos. 

O valor detectado só é totalmente rejeitado se o índice de clareza do sinal for inferior a $0.60$ (o equivalente a uma taxa de aperiodicidade de $0.40$). Isso assegura precisão de estúdio quando o som é limpo, sem punir ou ignorar a voz em desenvolvimento de um estudante que ainda carece de técnica vocal perfeita.

\subsection{Refinamento Parabólico}
A detecção discreta fica limitada a valores inteiros de amostras. Para frequências agudas, como um Dó Soprano ($C6 \approx 1047$ Hz), um erro de arredondamento de apenas $0.5$ amostra já causaria um desvio de afinação próximo a \textbf{20 cents}. Isso é inaceitável para treino auditivo, já que o ouvido humano percebe diferenças menores que 5 cents\footnote{A sensibilidade auditiva à variação de altura em ouvintes treinados pode ser inferior a 5 cents (aprox. 0.3\% da frequência) nas regiões centrais da fala e canto (Moore, B. C. J., 2012).}.

Para superar essa limitação física da taxa de amostragem sem precisar recorrer a um \textit{oversampling} custoso, aplica-se uma interpolação parabólica sobre o mínimo encontrado $\tau_{int}$. O período refinado $\tau_{fino}$ é calculado pelo vértice da parábola ajustada aos vizinhos do mínimo na função CMNDF:

\begin{equation}
    \tau_{fino} = \tau_{int} + \frac{d'[\tau_{int}-1] - d'[\tau_{int}+1]}{2(d'[\tau_{int}-1] - 2d'[\tau_{int}] + d'[\tau_{int}+1])}
\end{equation}

A frequência final é dada por $f_0 = f_s / \tau_{fino}$, garantindo precisão sub-cent sem adicionar latência.

\subsection{Cadeia de Filtragem e Estabilização Temporal}

A frequência bruta $\hat{f}_{raw}[n]$ obtida pelo YIN pode sofrer com dois tipos de instabilidade: erros impulsivos (saltos bruscos de oitava) e pequenas flutuações rápidas (\textit{jitter}). Para limpar isso sem comprometer a resposta em tempo real, o sistema usa uma cadeia de dois filtros: um não-linear seguido de um linear.

\subsubsection{Estágio 1: Filtragem de Mediana (Rank-Order Filter)}
Para eliminar erros grosseiros (\textit{outliers}) mantendo as transições reais entre as notas, o sinal passa primeiro por um filtro de mediana com janela $L=5$. Considerando o vetor de amostras recentes $W_n = [\hat{f}_{raw}[n], \dots, \hat{f}_{raw}[n-4]]$, a saída é:

\begin{equation}
    f_{med}[n] = \text{Mediana}(W_n)
\end{equation}

A escolha de uma janela ímpar de tamanho 5 permite ignorar até dois \textit{frames} errados consecutivos, introduzindo uma latência mínima de $(L-1)/2 = 2$ amostras, o que é essencial para a responsividade.

\subsubsection{Estágio 2: Suavização Exponencial (IIR)}
Em seguida, o sinal $f_{med}[n]$ passa por um filtro IIR de primeira ordem, que atua como um integrador com perdas. A equação é:

\begin{equation}
    f_{final}[n] = \alpha \cdot f_{med}[n] + (1 - \alpha) \cdot f_{final}[n-1]
\end{equation}

O valor de $\alpha$ foi definido em $0.50$. Esse é o ponto de equilíbrio: ele garante que a visualização da nota tenha firmeza e não fique tremendo na tela, mas sem deixar o sistema lento. Assim, ainda é possível capturar variações rápidas da voz.

\section{Mapeamento Musical (Conversão MIDI)}

Depois de obter a frequência fundamental $\hat{f}_0$ (em Hertz) pelo algoritmo YIN, o sistema converte esse valor para o espaço MIDI. Essa transformação é necessária para linearizar a percepção de altura. Na prática, isso garante que intervalos musicais iguais correspondam a distâncias numéricas constantes, não importa se a nota está numa oitava grave ou aguda.

\subsection{Modelagem Matemática}
O sistema usa como base o padrão internacional ISO 16:1975 \cite{iso16}, fixando o Lá central ($A_4$) em $\lambda = 440$ Hz. Isso garante compatibilidade com instrumentos ocidentais e serve de âncora para o índice MIDI $\mu = 69$. A função de mapeamento $\mathcal{M}: \mathbb{R}^+ \to \mathbb{R}$ é definida pela equação logarítmica:

\begin{equation}
    m = 69 + 12 \cdot \log_2 \left( \frac{\hat{f}_0}{440} \right)
\end{equation}

O resultado $m$ é um número real que guarda, ao mesmo tempo, a altura da nota e a precisão da afinação. Para saber qual é a nota musical exata $n$ na escala temperada, basta arredondar para o inteiro mais próximo:

\begin{equation}
    n = \lfloor m \rceil
\end{equation}

\subsection{Quantificação do Erro de Afinação (Biofeedback)}
A parte decimal que sobra define o desvio de afinação, chamado de $\delta = m - n$. No treino vocal, visualizar esse desvio é fundamental, já que a voz humana é um instrumento de altura contínua (não tem teclas fixas como um piano).

O sistema normaliza esse valor para a unidade de \textit{cents} (centésimos de semitom) usando a relação:
\begin{equation}
    c = 100 \cdot (m - n)
\end{equation}
O resultado $c$ fica no intervalo $[-50, +50]$. Esse dado alimenta a interface em tempo real, mostrando se o estudante está desafinando para baixo (bemol) ou para cima (sustenido). Assim, ele consegue corrigir a emissão vocal para chegar no centro tonal $n$, treinando o ouvido e o controle vocal.

\section{Análise de Estabilidade Vocal}

A seção anterior mostrou como obter a nota MIDI instantânea $m[t]$. Só que ter uma boa técnica vocal não é apenas acertar a afinação por um instante; é ter a capacidade de sustentar a emissão com consistência, evitando oscilações involuntárias (\textit{jitter} de altura) ou instabilidade por falta de apoio respiratório.

Para medir isso em tempo real, implementou-se uma métrica de estabilidade baseada na dispersão estatística de curto prazo.

\subsection{Modelagem Estatística (Janela Deslizante)}
O sistema mantém um \textit{buffer} circular com as últimas $K$ amostras de notas MIDI detectadas. O tamanho da janela foi definido em \textbf{$K=15$} quadros. Esse valor foi escolhido para oferecer um equilíbrio entre suavizar micro-flutuações irrelevantes e manter a reatividade para detectar tremolos reais.

Seja $\mathcal{M} = \{m_1, m_2, \dots, m_K\}$ o conjunto de valores na janela atual. O sistema calcula a média aritmética $\mu$ e a variância $\sigma^2$:

\begin{equation}
    \mu = \frac{1}{K} \sum_{i=1}^{K} m_i, \quad \sigma^2 = \frac{1}{K} \sum_{i=1}^{K} (m_i - \mu)^2
\end{equation}

O desvio padrão da amostra recente é dado por $\sigma = \sqrt{\sigma^2}$.

\subsection{Métrica de Pontuação Normalizada}
A pontuação de estabilidade $S \in [0, 1]$ é calculada normalizando o desvio padrão em relação a um limite de tolerância $\sigma_{lim}$.

No código, fixou-se $\sigma_{lim} = 0.5$ semitons. Esse valor representa o limite entre um vibrato artístico controlado e uma instabilidade indesejada. A função de pontuação final é:

\begin{equation}
    S = 1.0 - \min\left( \frac{\sigma}{\sigma_{lim}}, 1.0 \right)
\end{equation}

Dessa forma, uma emissão perfeitamente reta ("lisa") resulta em $S \approx 1.0$, enquanto oscilações que passam de meio semitom de amplitude derrubam a pontuação em direção a $0.0$. Isso fornece um \textit{biofeedback} visual imediato sobre o controle de ar e a tensão laríngea do estudante.

\section{Inferência de Tonalidade (Key Detection)}

Na seção anterior, o sistema obteve uma sequência de notas ao longo do tempo. Só que, na música ocidental, o significado de uma nota depende do contexto harmônico (a Tonalidade). Diferente da análise polifônica, onde os acordes deixam o tom óbvio, o Projeto Vaccai trabalha com sinal monofônico (apenas a melodia). Isso exige que o sistema descubra o contexto global a partir de uma linha melódica que, muitas vezes, contém erros de execução.

Para resolver essa ambiguidade das notas isoladas, utiliza-se um Modelo Oculto de Markov (HMM). A abordagem adapta o método de Müller \cite{muller_chap5} para priorizar a memória temporal e a estabilidade. A implementação detalhada está no \textbf{Algoritmo \ref{alg:key_detection}} (Apêndice A).

\subsection{Espaço de Estados e Observação}

O espaço de estados $S$ é formado pelas 24 tonalidades (12 maiores e 12 menores). Seja $q_t \in S$ a tonalidade oculta no instante $t$.
A variável observável $o_t$ não é a frequência absoluta, mas sim a \textbf{classe de altura} (ou \textit{chroma}), derivada da nota MIDI $m$ calculada na Seção 4:

\begin{equation}
    o_t = (m \mod 12) \in [0, 12)
\end{equation}

Essa operação de módulo torna o sistema "cego" à oitava. Ou seja, ele entende que um Dó central e um Dó agudo são a mesma classe de nota, permitindo detectar a tonalidade independente da voz ser grave (Baixo) ou aguda (Soprano).

\subsection{Probabilidade de Emissão (Modelo Fuzzy)}

A probabilidade de emissão $P(o_t | q_t = k)$ calcula a chance de se observar o chroma $o_t$ dado que a tonalidade atual é $k$. Para lidar com a afinação da voz humana (que tem vibratos e micro-desvios), o sistema não usa uma lógica binária ("pertence ou não pertence à escala"). Em vez disso, usa uma abordagem \textit{Fuzzy} Gaussiana.

Seja $d(o_t, S_k)$ a menor distância circular no relógio cromático entre a nota observada $o_t$ e as notas da escala $S_k$. A função é:

\begin{equation}
    E(o_t | k) = e^{-\frac{d(o_t, S_k)^2}{2\sigma^2}} + \epsilon
\end{equation}

\textbf{Justificativa dos Parâmetros:}
\begin{itemize}
    \item \textbf{$\epsilon = 0.01$:} Um valor mínimo de probabilidade para evitar erros numéricos (divisão por zero) em notas fora da escala.
    \item \textbf{$\sigma = 0.40$:} Esse valor define a tolerância. Foi estabelecido que uma nota com erro de afinação de $\pm 0.4$ semitons ainda conta com $\approx 60\%$ de relevância para a tonalidade correta. Isso permite que o sistema "perdoe" execuções imperfeitas, mantendo o contexto harmônico estável mesmo que o estudante desafine um pouco.
\end{itemize}

\subsection{Dinâmica de Transição e Inércia Tonal}

A matriz de transição $A$ define a probabilidade da música mudar de tom. Como modulações bruscas são raras em exercícios vocais, a matriz foi construída para favorecer a continuidade (auto-transição), criando uma espécie de "inércia".

As chances de transição para outros tons $A_{i,j}$ são baseadas na distância harmônica no \textbf{Círculo de Quintas}\footnote{O Círculo de Quintas organiza as 12 tonalidades. Nele, tonalidades vizinhas estão a um intervalo de quinta justa (7 semitons) e compartilham quase as mesmas notas.}. Isso permite modulações suaves (ex: Dó Maior $\to$ Sol Maior), mas penaliza saltos distantes e estranhos (ex: Dó Maior $\to$ F\# Maior).

\subsection{Atualização Causal Otimizada}

Em sistemas tradicionais de análise musical \textbf{offline}, algoritmos como o de Viterbi são frequentemente usados, pois analisam todo o áudio futuro para deduzir as escolhas corretas do passado. O motor DSP do Vaccai é estritamente de tempo real (causal), o que inviabiliza olhar para quadros futuros.

Para compensar isso e manter a estabilidade, o algoritmo considera a transição a partir da tonalidade mais provável do instante anterior ($k_{best}$) com um peso de inércia altíssimo. Seja $\mathbf{P}_t$ o vetor de probabilidades para as 24 tonalidades:

\begin{equation}
    \mathbf{P}_t[k] \propto E(o_t | k) \cdot \left( \alpha \cdot \mathbf{P}_{t-1}[k] + (1 - \alpha) \cdot A_{k_{best}, k} \right)
\end{equation}

Onde $\alpha = 0.9$ é o fator de inércia. Essa imposição causal (90\% de memória sobre 10\% de transição) reduz o custo de processamento a zero atraso perceptível e evita modulações errôneas originadas por notas isoladas fora do tom.


\clearpage

\appendix
\section{Implementação Algorítmica}

Abaixo apresentamos os algoritmos fundamentais que implementam o modelo matemático descrito. O Algoritmo \ref{alg:preprocess} detalha o condicionamento de sinal e a lógica de histerese. O Algoritmo \ref{alg:pitch} descreve a detecção de \textit{pitch} via YIN com pós-filtragem. Por fim, o Algoritmo \ref{alg:key_detection} descreve a inferência harmônica Fuzzy.

\begin{algorithm}[H] 
\caption{Pré-processamento: Noise Gate com Histerese e ZCR}
\label{alg:preprocess}
\begin{algorithmic}[1]
\Require Buffer de áudio bruto $x$ de tamanho $N$, Estado anterior $S_{prev}$
\Ensure Buffer centralizado $x'$, Novo Estado $S$

\State $\mu \gets \frac{1}{N} \sum x[i]$ \Comment{Média DC}
\State $E \gets 0, \quad \text{crossings} \gets 0$

\For{$i \gets 0$ \textbf{to} $N-1$}
    \State $x'[i] \gets x[i] - \mu$ \Comment{Remoção de Offset}
    \State $E \gets E + (x'[i])^2$
    \If{$i > 0$ \textbf{and} $\text{sgn}(x'[i]) \neq \text{sgn}(x'[i-1])$}
        \State $\text{crossings} \gets \text{crossings} + 1$
    \EndIf
\EndFor

\State $rms \gets \sqrt{E / N}$
\State $zcr \gets \text{crossings} / (N-1)$

\State \textbf{Lógica de Histerese Dupla (Schmitt Trigger)}
\If{$rms > 0.003$ \textbf{and} $zcr < 0.20$}
    \State $S \gets \textbf{True}$ \Comment{Abre o Gate (Voz Limpa Detectada)}
\ElsIf{$rms < 0.001$ \textbf{or} $zcr > 0.30$}
    \State $S \gets \textbf{False}$ \Comment{Fecha o Gate (Silêncio ou Ruído de Sopro)}
\Else
    \State $S \gets S_{prev}$ \Comment{Mantém estado na Zona Morta}
\EndIf

\If{$S$ \textbf{is} \textbf{True}}
    \State \Return \textbf{True}, $x'$ 
\Else
    \State \Return \textbf{False}, $\emptyset$ 
\EndIf
\end{algorithmic}
\end{algorithm}

\begin{algorithm}[H]
\caption{Estimativa de $f_0$ via Algoritmo YIN com Filtragem Híbrida}
\label{alg:pitch}
\begin{algorithmic}[1]
\Require Sinal $x'$, Taxa $f_s$, Histórico $f_{prev}$
\Ensure Frequência Estabilizada $f_{final}$

\State $\tau_{min} \gets f_s / 1000, \quad \tau_{max} \gets f_s / 50$
\State Inicializar vetor de diferença $d[\tau]$ e CMNDF $d'[\tau]$

\State \textbf{Passo 1: Diferença Quadrática e CMNDF}
\State $d'[0] \gets 1, \quad \text{running\_sum} \gets 0$
\For{$\tau \gets 1$ \textbf{to} $\tau_{max}$}
    \State $d[\tau] \gets \sum (x'[i] - x'[i+\tau])^2$
    \State $\text{running\_sum} \gets \text{running\_sum} + d[\tau]$
    \State $d'[\tau] \gets d[\tau] \cdot (\tau / \text{running\_sum})$
\EndFor

\State \textbf{Passo 2: Busca Híbrida por Vale (Limiar + Fallback)}
\State $\hat{\tau} \gets -1, \quad \text{min\_val} \gets \infty$
\For{$\tau \in [\tau_{min}, \tau_{max}]$}
    \If{$d'[\tau] < 0.15$} \Comment{Limiar Rigoroso}
        \While{$d'[\tau+1] < d'[\tau]$} \Comment{Busca mínimo local}
            \State $\tau \gets \tau + 1$
        \EndWhile
        \State $\hat{\tau} \gets \tau$
        \State $\text{min\_val} \gets d'[\hat{\tau}]$
        \State \textbf{Break}
    \EndIf
    \If{$d'[\tau] < \text{min\_val}$}
        \State $\text{min\_val} \gets d'[\tau]$ \Comment{Grava para o Fallback}
        \State $\text{global\_min\_tau} \gets \tau$
    \EndIf
\EndFor

\If{$\hat{\tau} == -1$}
    \State $\hat{\tau} \gets \text{global\_min\_tau}$ \Comment{Aplica Fallback se não bater 0.15}
\EndIf

\State \textbf{Passo 3: Validação por Claridade}
\State $clarity \gets 1.0 - \text{min\_val}$
\If{$clarity < 0.60$} \Comment{Corte de pureza final (0.40 aperiodicidade)}
    \State \Return $0$
\EndIf

\State $f_{raw} \gets f_s / \text{ParabolicInterpolation}(\hat{\tau})$

\State \textbf{Passo 4: Filtragem Híbrida (Mediana + IIR)}
\State $f_{med} \gets \text{MedianFilter}(f_{raw}, \text{buffer\_size}=5)$
\State $f_{final} \gets 0.5 \cdot f_{med} + 0.5 \cdot f_{prev}$ \Comment{Suavização IIR}

\State \Return $f_{final}$
\end{algorithmic}
\end{algorithm}

\begin{algorithm}[H]
\caption{Detecção de Tonalidade Causal (HMM Gaussiano Fuzzy)}
\label{alg:key_detection}
\begin{algorithmic}[1]
\Require Nota MIDI $n$, Vetor de Probabilidades $\mathbf{P}_{prev}$
\Ensure Tonalidade Vencedora $K_{win}$

\State $c \gets n \pmod{12}$ \Comment{Nota cromática contínua}
\State $\mathbf{P}_{new} \gets \text{zeros}(24)$
\State $k_{best\_prev} \gets \text{argmax}(\mathbf{P}_{prev})$
\State $\sigma \gets 0.40$ \Comment{Parâmetro de tolerância}

\For{$k \gets 0$ \textbf{to} $23$}
    \State \textbf{Passo 1: Emissão Gaussiana}
    \State $d_{min} \gets \infty$
    \For{$degree \in \text{ScaleDegrees}[k]$}
        \State $d \gets \text{CircularDistance}(c, degree)$
        \State $d_{min} \gets \min(d_{min}, d)$
    \EndFor
    \State $E \gets e^{-(d_{min}^2 / 2\sigma^2)} + 0.01$

    \State \textbf{Passo 2: Transição e Inércia Causal}
    \State $T \gets \text{TransitionMatrix}[k_{best\_prev}][k]$
    \State $\mathbf{P}_{new}[k] \gets E \cdot (0.9 \cdot \mathbf{P}_{prev}[k] + 0.1 \cdot T)$
\EndFor

\State \textbf{Passo 3: Normalização}
\State $S \gets \sum \mathbf{P}_{new}$
\If{$S > 0$}
    \State $\mathbf{P}_{prev} \gets \mathbf{P}_{new} / S$ \Comment{Atualiza e normaliza}
\Else
    \State $\mathbf{P}_{prev} \gets \text{Uniform}(24)$ \Comment{Fallback}
\EndIf

\State $K_{win} \gets \text{argmax}(\mathbf{P}_{prev})$

\State \Return $K_{win}$
\end{algorithmic}
\end{algorithm}

\begin{thebibliography}{9}

% Seção 3: YIN
\bibitem{yin2002}
De Cheveigné, Alain; Kawahara, Hideki.
``YIN, a fundamental frequency estimator for speech and music''.
\textit{The Journal of the Acoustical Society of America}, vol. 111, no. 4, pp. 1917-1930.
2002.

% Seção 3: Nota de Rodapé (JND)
\bibitem{moore_hearing}
Moore, Brian C. J.
\textit{An Introduction to the Psychology of Hearing}.
6th Edition. Emerald Group Publishing Limited, 2012.

% Seção 5: ISO MIDI
\bibitem{iso16}
International Organization for Standardization.
\textit{ISO 16:1975 Acoustics — Standard tuning frequency (Standard musical pitch)}.
1975.

% Seção 6: HMM
\bibitem{muller_chap5}
Müller, Meinard.
\textit{Fundamentals of Music Processing: Audio, Analysis, Algorithms, Applications}.
Capítulo 5: Harmony Analysis.
Springer, 2015.
ISBN 978-3-319-21944-8.

\end{thebibliography}

\end{document}